\section{Related Work}\label{sec:related}

LRL combines ideas from several lines of work.
For each line we say what LRL shares with it and where it differs.
Section~\ref{sec:eval:compare} compares LRL with Lean~4, Rust, Idris~2 and Racket by running programs; here we position the design.

\subsection{Linear and dependent type theories}\label{sec:related:lindep}

Linear logic~\cite{girard1987linear} and Wadler's proposal to use linear types for safe in-place update~\cite{wadler1990} started the study of resource-sensitive types; combining them with dependent types is a line of work of its own.
LLF~\cite{cervesato2002llf} adds linear implication to the LF logical framework but lets types depend only on unrestricted terms.
Later calculi keep the same separation between a dependent, unrestricted world and a linear world~\cite{krishnaswami2015integrating}, and such systems have been given a categorical semantics~\cite{vakar2014syntax,vakar2015categorical}, including one for quantum circuits~\cite{fu2022linear}.
Quantitative Type Theory (QTT)~\cite{mcbride2016,atkey2018} annotates every binder with a usage count drawn from a semiring and counts uses inside types as zero, so that linear variables may appear in types; Idris~2 implements QTT~\cite{brady2021}.
Graded dependent type theories also track how variables are used in types~\cite{moon2021grtt} and give usage an operational meaning~\cite{choudhury2021grad}, and one with erasure has been formalised in Agda~\cite{abel2023graded}.
A two-level design~\cite{fu2026twolevel} separates a logical level from a linear program level whose proofs and types are erasable; ATS~\cite{xi2007ats,zhu2005stateful,chen2005combining} combines dependent types, linear proofs about memory, and proofs written as total functions that are erased before execution.

LRL differs from these systems in two ways.
First, its annotations describe something else.
QTT quantities, and the multiplicities of Linear Haskell~\cite{bernardy2018}, describe how a function uses its \emph{argument}.
LRL's function kinds describe how a closure uses its \emph{captured environment}, as Rust's \lstinline|Fn|/\lstinline|FnMut|/\lstinline|FnOnce| traits do (\S\ref{sec:core:kinds}).
The closest published analogue is the affinity mode of Oxidizing OCaml~\cite{lorenzen2024oxidizing}, in which a closure that captures a unique value can be called only once.
Second, LRL keeps typing and ownership as two separate judgments in its kernel---occurrences of an owned variable inside a type are erased uses and impose no constraint---and checks borrows with a third, flow-sensitive analysis on its intermediate representation.
Work on graded types takes the other route, folding uniqueness, ownership and borrowing into one graded type system~\cite{marshall2022,marshall2024fractional}, and Pure Borrow~\cite{matsushita2026pureborrow} brings Rust-style borrowing to Linear Haskell.
Section~\ref{sec:tensions} discusses what our separation costs.
Unlike QTT and graded type theories, LRL has no soundness proof for its full dependent calculus; our machine-checked results cover a non-dependent core (\S\ref{sec:own:meta}).

\subsection{Languages with linearity, uniqueness, or ownership of state}\label{sec:related:lang}

Clean's uniqueness types~\cite{barendsen1993conventional,barendsen1996uniqueness}, Alms~\cite{tov2011practical}, Linear Haskell~\cite{bernardy2018}, Granule~\cite{orchard2019granule} and Cogent~\cite{oconnor2016cogent} are implemented languages that statically restrict duplication; Mezzo~\cite{pottier2013mezzo,balabonski2016mezzo} tracks duplicable and affine permissions that change from one program point to the next.
Vault~\cite{deline2001vault} enforces resource protocols in low-level code, and alias types~\cite{smith2000alias} track the shape of the store and the aliasing of pointers, so that low-level code can deallocate and reuse memory safely.
Session types~\cite{honda1998session}, their reading in linear logic~\cite{caires2010session,wadler2012sessions}, and their dependent extension~\cite{toninho2011dependent} are the standard account of protocols such as the channel of \S\ref{sec:tour}.
Systems that reason about mutable state with proofs inside a dependently typed language---Hoare Type Theory and Ynot~\cite{nanevski2006htt,nanevski2008htt,nanevski2008ynot}, Low*~\cite{protzenko2017lowstar} with explicit memory footprints, and the separation logics Steel and PulseCore in F*~\cite{swamy2016fstar,swamy2020steelcore,fromherz2021steel,ebner2025pulsecore}---reason about ownership of mutable state with proofs rather than with an affine type discipline.
They combine dependent types with reasoning about ownership of mutable state far beyond what LRL offers; LRL, which has no mutable state at all (\S\ref{sec:limits}), is much less expressive in this respect.

\subsection{Refinement types and verification for Rust}\label{sec:related:rust}

RustBelt~\cite{jung2018rustbelt} proves a substantial fragment of Rust sound in the Iris logic; Oxide~\cite{weiss2021oxide} formalises borrow checking; Stacked and Tree Borrows~\cite{jung2020stacked,villani2025treeborrows} are aliasing models for unsafe code; LRL's MIR borrow checker follows non-lexical lifetimes~\cite{rfc2094}, computing each region as a set of program points; Polonius~\cite{polonius2018} reformulates the same check in terms of loans.
A second line adds specifications and proofs to Rust programs: Prusti~\cite{astrauskas2019prusti}, Creusot~\cite{denis2022creusot}, Verus~\cite{lattuada2023verus,lattuada2024verussosp}, RustHorn~\cite{matsushita2020rusthorn}, Aeneas~\cite{ho2022aeneas}, RefinedRust~\cite{gaher2024refinedrust} and the Kani model checker~\cite{vanhattum2022kani}; RefinedC~\cite{sammler2021refinedc} does the same for C.
Flux~\cite{lehmann2023flux} adds refinement types to Rust, checks them on rustc's MIR and indexes vectors by their length; Thrust~\cite{ogawa2025thrust} adds a dependent refinement type system that uses prophecy variables for mutable borrows.
These tools can state and prove properties like those of our case studies, but they are separate tools layered on a mature language whose compiler does not check the proofs.
LRL has one kernel that checks both programs and proofs, with dependent types in the style of the Calculus of Inductive Constructions rather than refinements, and none of Rust's maturity.

\subsection{Dependent types and hygienic macros}\label{sec:related:macros}

Hygienic macro expansion goes back to the 1980s and early 1990s~\cite{kohlbecker1986,dybvig1993}; scope sets~\cite{flatt2016scopes} underlie hygiene in Racket and inspired the simpler macro scopes of Lean~4~\cite{ullrich2022}.
Turnstile~\cite{chang2017turnstile} and Turnstile+~\cite{chang2020dependent} reverse the usual relationship between macros and types: the type checker itself is written as macros, which makes it possible to build dependently typed languages, up to the Cur proof assistant, inside Racket.
Typed Racket~\cite{tobinhochstadt2008typedscheme,tobinhochstadt2011languages} and its refinement types~\cite{kent2016occurrence} are other typed languages built from macros, and macro-extensible domain-specific languages can be built in the same way~\cite{ballantyne2020macros}.
Template Haskell~\cite{sheard2002th}, Scala~3 macros~\cite{stucki2018scala,stucki2021multistage}, Lean~4 elaborators~\cite{ullrich2022}, Idris elaborator reflection~\cite{christiansen2016elab} (also available in Idris~2) and Agda's reflection~\cite{vanderwalt2013agda} let metaprograms inspect types; Rust's \lstinline|macro_rules!| has what the Rust reference calls mixed-site hygiene~\cite{rustmacros}.
LRL takes the most restrictive position in this space: macros are templates with no conditionals, no pattern matching on arguments and no compile-time evaluation (\S\ref{sec:macros:alg}); they cannot inspect types or see ownership information, and their output enters the pipeline before elaboration.
This gives up expressiveness compared with every system above, in exchange for a guarantee that is easy to state and to test: a macro can only produce programs that a user could have written by hand, up to renaming of bound variables; those programs pass the same checks as hand-written ones, and in addition a macro may not produce axioms, unsafe definitions, \lstinline|eval| or classical imports (\S\ref{sec:macros:boundary}).
Turnstile, where type checking happens during expansion, is the opposite design point.

\subsection{Ownership types in the object-oriented tradition}\label{sec:related:oo}

Object-oriented languages studied the control of aliasing long before Rust: Islands~\cite{hogg1991islands}, Balloons~\cite{almeida1997balloon}, flexible alias protection~\cite{noble1998flexible} and ownership types~\cite{clarke1998ownership}, surveyed in~\cite{clarke2013survey}.
Later work added uniqueness and borrowing: alias burying~\cite{boyland2001burying}, which checks uniqueness with a liveness analysis and so anticipates non-lexical lifetimes; capabilities for sharing~\cite{boyland2001capabilities}; external uniqueness~\cite{clarke2003external}; AliasJava~\cite{aldrich2002alias}; ownership domains~\cite{aldrich2004domains}; ownership types for freedom from data races and deadlocks~\cite{boyapati2002ownership}; capabilities for uniqueness and borrowing~\cite{haller2010capabilities}; reference immutability~\cite{gordon2012uniqueness}; the reference capabilities of Pony~\cite{clebsch2015deny}; and reference capabilities for data-race-free concurrency control~\cite{castegren2016reference}.
Universe types use ownership to support modular verification~\cite{muller2002modular,dietl2005universes,dietl2007generic}; existential owners~\cite{wrigstad2007existential} and loose ownership domains~\cite{schaefer2007loose} make ownership types more expressive; and Joe3~\cite{ostlund2008ownership} combines ownership with uniqueness and immutability.
Recent work in this journal studies mutable value semantics as an alternative to references~\cite{racordon2022mvs} and a capability-based intermediate language that guarantees uniqueness and immutability~\cite{racordon2022lingua}.
In this paper ``ownership'' means Rust's discipline of moves and borrows---every value has one owner, a move transfers it, and references are checked against regions---and not owner-parameterised types, which constrain the shape of an object graph.
LRL has no objects, no mutable heap and no owner parameters.
Its closures come closest: as \S\ref{sec:core:kinds} explains, a closure is an object whose fields are its captures, and its kind says whether its single method takes the receiver by shared reference, by mutable reference or by value.
The two traditions meet in uniqueness and borrowing~\cite{boyland2001burying,haller2010capabilities}: in LRL a value whose type is not Copy, including every closure that is not erased (\S\ref{sec:core:copy}), has a single owner and can be borrowed, much as a unique reference can be. A closure's kind is a separate property: \lstinline|FnOnce| limits how often a closure may be called, which Oxidizing OCaml calls affinity and distinguishes from uniqueness~\cite{lorenzen2024oxidizing}.

\subsection{In-place update in functional and dependently typed languages}\label{sec:related:inplace}

Static guarantees of destructive update go back to linear types~\cite{wadler1990}, Clean's uniqueness types~\cite{barendsen1996uniqueness} and Hofmann's type system for in-place update~\cite{hofmann2000inplace}.
Lean~4~\cite{lean42021} and Koka~\cite{leijen2014koka} instead use precise reference counting to update unshared values in place at run time~\cite{ullrich2019beans,reinking2021perceus,lorenzen2022frame}---a style that Perceus calls ``functional but in-place'' (FBIP). FP\textsuperscript{2}~\cite{lorenzen2023fp2} checks statically that a function can run fully in place, without allocating memory, provided its arguments are not shared; its implementation in Koka checks that condition at run time.
Idris~2 offers linear arrays updated in place~\cite{brady2021}.
We observed the run-time behaviour of Lean~4 directly (\S\ref{sec:eval:compare}): it reversed an unshared array without copying it and copied a shared one.
LRL performs no in-place update at all: its values are immutable and its backends copy them.
Its ownership discipline would make a static reuse analysis possible, but we have not built one, and we make no performance claim based on ownership.

\subsection{Trusted kernels}\label{sec:related:kernels}

LRL follows proof assistants such as Coq~\cite{bertot2004coqart} and Lean~\cite{lean2015,lean42021} in routing every definition through a separate checker whose correctness suffices for the logical soundness of what it accepts.
Verified compilers such as CompCert~\cite{leroy2009compcert} and CakeML~\cite{kumar2014cakeml} go much further and prove, in a proof assistant, that their compilation passes preserve the meaning of programs; LRL proves nothing about its compiler.
The trust question specific to LRL is which resource properties lie inside the kernel's guarantee and which do not; \S\ref{sec:own:trust} answers it.
